\subsection{Related inverse problems}
\label{sec:theorem-comparison}
The closest comparisons concern geometric factorization and additive
deconvolution with unknown noise. We describe their data and uniqueness
classes explicitly. This also distinguishes the identification supplied
by the collision fields from the classical inversions used after the
geometric factors have been separated.

For planar convex bodies $C,A$, write
\[
 g_{C,A}(x)=|C\cap(A+x)|
          =(\mathbf1_C*\mathbf1_{-A})(x).
\]
Averkov and Bianchi \cite[Theorem~1.1]{AverkovBianchi} prove that
$g_{C,C}$ determines every planar convex body up to translation and
reflection. The corresponding problem with two unknown bodies has
also been solved in substantial classes. Bianchi
\cite[Theorem~1.1]{BianchiCross} determines a pair of planar polygons
from its cross covariogram, apart from specified parallelogram
exceptions and the unavoidable common translations and reflected
interchanges $(C,A)\mapsto(-A,-C)$. For planar bodies of class
$C^8_+$, meaning $C^8$ boundaries with positive curvature,
\cite[Theorem~6.3]{BianchiLaplace} proves uniqueness up to precisely
these two equivalences. Thus simultaneous recovery of two convex
factors is an established geometric inverse problem.

When the launch law is uniform on a full-dimensional convex body $A$,
our isolated occupation field is $|A|^{-1}g_{C,A}$; the normalization
is part of the comparison. Theorem~\ref{thm:two-field-rigidity} starts
with the ordered collision fields, permits an arbitrary compact law
with convex support, and imposes strict convexity only on the obstacle.
The recovered collision supports provide the additional factorization
identity. Its full fiber is common translation, without a reflected
interchange. The geometric conclusions also extend to separated,
locally finite obstacle configurations. On the statistical side,
Bianchi, Gardner and Kiderlen \cite[Theorems~4.10 and~6.4]{BGK} give
strong Hausdorff consistency for reconstruction from noisy covariogram
grids, using a preliminary difference-body reconstruction and their
convex-body uniqueness and noise assumptions. The finite
experiment below acquires Bernoulli collision bits and controls the
resulting geometric and probability errors in that experiment.

Unknown-probe morphology gives another geometric comparison.
Villarrubia's blind tip reconstruction
\cite[Section~5.1, equations~(14)--(15)]{Villarrubia} successively
intersects dilation constraints, starting from a containing tip bound.
The limiting reconstruction is the largest tip consistent with those
constraints; additional images can be incorporated
\cite[Section~7.3.3]{Villarrubia}. This identifies a maximal consistent
probe, with exact recovery governed by the available contacts.
In the present problem the three matched supports determine both
geometric factors, and the occupation values subsequently determine
the probability carried by the footprint.

For passive additive observations $Y=X+\epsilon$, Gassiat, Le Corff
and Leh\'ericy \cite[Theorem~2.1]{GLL} identify both unknown laws,
up to opposite translations of the summands, when the noise has two
independent coordinate blocks and the signal has an entire Laplace
transform of order less than two satisfying their condition (H2).
Writing the entire characteristic transform as $\Phi$, this condition
requires
\[
 \Phi(z_1,\,\cdot\,)\not\equiv0
 \quad\hbox{and}\quad
 \Phi(\,\cdot\,,z_2)\not\equiv0
 \qquad\hbox{for every complex }z_1,z_2.
\]
Their theorem permits singular laws and zeros of the noise transform.
Capitao-Miniconi, Gassiat and Leh\'ericy
\cite[Corollary~2.5]{CGLSupport} give geometric sufficient conditions
for this hypothesis, including compact strictly convex supports.
Their Theorem~5.1 gives $W_p$ risk for the compact signal law of order
$\log\log n/\log n$ on the stated compact-support classes, which
include a quantitative contrast condition in addition to the exact
identification assumptions.

There is a direct overlap with our occupation reduction. If $U_C$ is
uniform on the recovered component $C$ and independent of the launch
displacement $Z$, the normalized occupation density is the law of
$U_C-Z$. When $Z$ has independent coordinate blocks, the preceding
identification theorem applies, since $C$ is strictly convex. Our
launch law need not have this independence. The two active collision
fields supply the signed support identity that separates the factors
for a general planar launch law. Neither the passive additive samples
nor their statistical rates are assumed available in the finite
collision experiment.

The support-measure step has a precise classical antecedent.
Martinez-Maure \cite[Definition~4.4 and Theorem~4.6]{MartinezMaure}
associates a signed length measure to a difference of support
functions and determines that difference, modulo translations, from
its zero-first-moment measure. Such a difference alone does not
determine its two convex summands. Here the observable identity yields
$h_C-h_{L_C}$. Strict convexity makes $S_C$ nonatomic, whereas the
segment contributes exactly two atoms. The Jordan decomposition
therefore separates the summands under the hypotheses of our exact
theorem, including nonsmooth strictly convex obstacles.

Compact support and Fourier zeros also have a substantial
deconvolution literature. Meister \cite[Section~2 and Theorem~1]{MeisterCompact}
recovers a compactly supported univariate density using the known
error transform only near the origin. Delaigle and Meister
\cite[Theorem~4.1]{DelaigleMeister} obtain statistical bounds for
univariate deconvolution with specified oscillating error transforms
and one-sided signal support. In our exact planar inverse the compact
body indicator has an entire Fourier transform, nonzero at zero.
Its real nonzero set is dense, so division there and continuity
determine the launch transform everywhere. This argument allows
nonisolated planar zero sets. The finite result instead uses mixed
moments with explicit conditioning for the estimated body factor.

Moment control in transportation distance is classical; see
Rigollet and Weed \cite[Section~2.3, Theorem~4 and Corollary~1]{RigolletWeed}
for quantitative univariate bounds. The proof below supplies the
bivariate approximation and convolution estimates needed after the
geometry has been estimated, together with their acquisition from
finite collision bits. The prediction conclusions then use standard
variation bounds \cite{AFP,Galerne}. The resulting identification
chain and finite resource bounds concern the same two prescribed
positive-length commands throughout.
